\chapter*{Appendix 13: Scenario Nu (Fashion)}

\section*{1. The Legacy Dependency (Technical \\\& Cultural)}
The legacy textile and fashion industry is built on hyper-extractive fast fashion, globalized sweatshop labor, and synthetic petroleum-based materials. The dependency is largely cultural (status signaling through legacy brands) and technical (reliance on globalized industrial looms and fiat-driven supply chains).

\section*{2. The Parasitic Bridge (Technical \\\& Legal)}
Layer 7 interfaces with local, legacy textile producers through "Micro-Manufacturing Guild DAOs." These DAOs use fiat wrappers to lease time on legacy industrial equipment during off-peak hours (Grid Leeching). Culturally, the bridge subverts legacy fashion by embedding sovereign cryptographic identifiers (NFC chips/Layer 5 tokens) into locally produced garments, repurposing the memetic desire for status into a proof-of-sovereignty mechanism.

\*{3. The Decoupling Trigger}
The reliance on legacy industrial leasing is decoupled when the sovereign network achieves sufficient density of decentralized, open-source automated looms and bio-textile synthesis (e.g., mycelium leather) to meet local demand.

\section*{4. Orchestration \\\& Ethical Mechanics}
The orchestration involves smart contracts that transparently trace every thread and labor hour, diametrically opposing the opacity of legacy fashion. Ethically, the bridge aims to absorb the skills of legacy garment workers into the sovereign mesh, compensating them in Layer 3 assets that hold higher purchasing power in the local network than their previous fiat wages, thus starving the legacy sweatshop system of labor.
